\section{Benchmark and governance}
\label{sec:benchmark}

\subsection{Governed trace contract}

Every trajectory in the benchmark is governed by a typed event contract
that extends the underlying \textsc{Autoresearch} harness~\citep{autoresearch2026}.
The contract specifies fourteen event types—including proposal, action
request, action result, state snapshot, state difference, evidence,
verifier request, verifier result, fault injection, belief, recovery
action, human escalation, terminal outcome, and adjudication—each with
mandatory fields and parent-event links for causal traceability. A stable
envelope provides run identity, monotonic step index, wall-clock timestamp,
actor type, artifact references, and a visibility class that controls
whether the payload is policy-visible, hidden fault truth, sealed test
outcome, training-reward-only, or adjudication-only.

The state-difference contract records the target resource (file, process,
dataset, metric, or environment), before/after content-addressed snapshots
or explicit missing-state markers, the semantic type of change, the
postcondition expected before execution, the observed change, whether an
interrupted or partial execution was independently reconciled, and the
observer's confidence. Absence—timeout, missing observation, failed
verifier, unexecuted action—is typed rather than collapsed into a single
null value.

An immutable run manifest is hashed before execution. It records task
identity, model version, environment lock, resource budget, experimental
assignment, a sealed fault-manifest hash, and adjudicator versions. A
manifest edited after execution is invalid.

\subsection{Causal fault model}

We organize controlled faults into seven mechanism families, each with a
declared causal parent, affected component, and injection trigger:

\begin{enumerate}
    \item \textbf{Model and planning:} incorrect hypothesis, invalid task
          decomposition, or hallucinated dependency. The error originates in
          the reasoning step before action.
    \item \textbf{Action and tool:} malformed tool call, wrong target, or
          invalid parameter. Intent differs from executed effect.
    \item \textbf{Runtime and environment:} timeout after a side effect,
          stale process state, or partial write. Reconciliation is required
          to recover the true state.
    \item \textbf{Memory and context:} stale fact, retrieval corruption,
          context truncation, or provenance loss. A memory event changes the
          agent's plan before it reaches the tool layer.
    \item \textbf{Evidence and statistics:} data leakage, incorrect metric,
          multiple-testing violation, or invalid uncertainty estimate. The
          scientific claim is unsupported despite task completion.
    \item \textbf{Verifier:} flaky test, shared implementation bug between
          agent and verifier, incomplete oracle, or reward leakage.
          Acceptance is not independent of the agent's own reasoning.
    \item \textbf{Coordination:} conflicting worker instructions, stale
          ownership, or duplicated action in a multi-agent configuration.
          Failure propagates across agent boundaries.
\end{enumerate}

Each injected fault carries a deterministic seed, severity (low, medium, or
high), persistence (transient or persistent), an observability class, and a
list of expected immediate and downstream effects. Every fault assignment
has a clean counterfactual—the same task executed without the fault—to
enable within-task causal estimation. The fault truth is sealed in a
manifest that the evaluated policy cannot read.

\subsection{Workload ladder}

Tasks are drawn from four workload levels of increasing complexity:

\begin{description}
    \item[L0---Repository repair:] bounded code-editing tasks with precise
          state-difference labels and fast iteration. Used for injector
          debugging and fidelity measurement.
    \item[L1---Reproducible ML experiments:] training, evaluation, and
          hyperparameter tasks where data leakage, configuration drift, and
          metric misuse are observable.
    \item[L2---Signal-detection workflows:] tasks with known optimal
          baselines and controlled noise geometry, enabling measurement of
          statistical and evidence-level faults.
    \item[L3---Nonstationary-noise experiments:] tasks with distribution
          shift, stale evidence, and model adaptation under changing
          conditions.
\end{description}

L0--L3 support causal injection and mechanism analysis. An L4 real-data
transfer arm, using naturally occurring failures without controlled fault
labels, is included for external validity but does not support causal claims
about specific mechanisms.

\subsection{Split regime and leakage controls}

Six split regimes are predeclared and frozen before confirmatory collection.
Development splits consist of pilot-only task templates and injectors used
for instrument repair. The confirmatory evaluation uses: (i)~unseen task
templates, where entire task generators and their instances are withheld
from training; (ii)~unseen injector implementations, where a known causal
mechanism is implemented through a new code path; (iii)~unseen fault
families, where one mechanism family is absent from model development
(attribution for this regime is evaluated as ``unknown'' rather than forced
into a known class); (iv)~unseen compositions, where individually familiar
faults appear in new pairs or sequences; and (v)~an external L4 split with
naturally occurring failures.

Task templates, repositories, datasets, and near-duplicate solutions are
grouped before any split assignment. Text and code similarity screens
identify borderline clusters, which are manually audited. No task,
repository, or injector implementation may appear in more than one
confirmatory split.

\subsection{Adjudication and evidence separation}

Adjudication separates execution success and scientific validity. A
trajectory may complete execution correctly but still produce a
scientifically invalid claim (e.g., a correct tool call on incorrectly
preprocessed data). The adjudication record captures deterministic checks,
scientific postconditions, verifier versions, disagreements across
adjudicators, human review decisions, and the final status from the
vocabulary \{accepted, rejected, ambiguous, insufficient\_evidence,
verifier\_failure\}. Ambiguous cases retain their ambiguity rather than
being coerced into a majority class.

The training reward channel, agent self-report, and final adjudication are
maintained as three separate evidence streams. An outcome that satisfies the
training verifier but fails independent scientific adjudication is recorded
as verifier-exploitable rather than silently accepted.

\subsection{Artifact access, licensing, and maintenance}

The benchmark is released with an executable trace schema and validator,
deterministic fault injectors and clean counterfactuals, frozen task, fault,
model, and split manifests, benchmark traces with provenance and licensing
records, evaluation code for survival, warning, calibration, attribution,
and recoverability, and a documented maintenance plan. The public release
includes a de-identified trace sample and the full evaluation harness;
sensitive or access-controlled traces are referenced through archival
identifiers with documented retrieval instructions.

\begin{figure}[t]
\centering
\fbox{\parbox{0.9\textwidth}{\centering
\textbf{Architecture diagram:} Governed harness with evidence channels,
causal fault intervention model, and split geometry.\\
\textit{[Replace with accessible text-source diagram using Mermaid or TikZ; see source in figures/]}}}
\caption{System architecture showing the governed runtime, causal fault injection,
split isolation, and the relationship between evidence channels and
adjudication.}
\label{fig:benchmark_architecture}
\end{figure}
